\section{Procedure and commitments} \label{app:procedure}

\begin{enumerate}
\item \textbf{Order.} No Sector 110 product of any kind (light curves, pixels, full-frame images, SPOC data-validation
  reports, threshold-crossing-event or TESS Object of Interest lists for this star) is opened before the primary run.
  The primary run is made within 30 days of the release of the primary product (Appendix~\ref{app:data}). A later run is
  still the primary run, and its delay is reported.
\item \textbf{Refusals.} Every registered run (photometry, variants, localisation) refuses to start, before reading any
  data, in three cases: the installed package versions, the machine type or the operating system differ from
  \code{requirements.txt}; the input is not Sector 110 of TIC 4206066; or the file name is not the registered name of
  the product. If MAST distributes a product under another name, the file is renamed to the registered pattern of
  \code{s110\_final.py}, and the original name and its sha256 are recorded. A refused run writes nothing and is not a
  run. The exact commands, run from the root of the deposit, are listed in
  Appendix~\ref{app:commands}.
\item \textbf{First run and posting.} The first primary run that writes its output file is the primary run. That file
  (which contains the input, code and dependency checksums and the package versions) and its sha256 are posted as a new
  version of this Zenodo record within 7 days of the run. Only then are the variants and the localisation run, within
  14 days of posting the primary result. Their output files and the summary file are posted in the same way within 7
  days of being made.
\item \textbf{Failures and repeats.} If a script writes its output file, that file defines its outcomes, even if a later
  step fails. A run that ends before writing its output for any reason other than a refusal or a failure to read the
  file (for example an interruption) is repeated once, unchanged. If it fails again, the outcome is O4 for both signals
  (photometry) or ``localisation not possible''. For a reading failure the unedited error message of the script is saved
  and its checksum recorded. A code edit is allowed only to read the file (column or extension names, data type), and
  the edited script keeps its registered file name (Appendix~\ref{app:p6}). It may not change any constant, window,
  mask, seed, threshold or rule. Before use, the edited script must reproduce exactly the \code{results} block of every
  photometric \code{dryrun/*.json}, and the \code{stats} block and every \code{injections.sets.*.trials} list of every
  localisation \code{dryrun/loc\_*.json} (run times differ and are not compared). The diff and this check are posted, as
  in point 3, before the Sector 110 run. Variants are never edited. Any other failure of the localisation gives
  ``localisation not possible''.
\item \textbf{Full reporting.} All outcomes are reported, whatever they are, for A and B, with all variants, the
  transits that were not tested and the reasons, and the localisation.
\item \textbf{No changes.} Nothing in this deposit is changed after publication. Any later idea is a new, separately
  labelled analysis.
\item \textbf{Public checks.} Readers can check every posted output file: its \code{script\_sha256},
  \code{deps\_sha256}, \code{ephemeris\_sha256} and \code{expected\_depths\_sha256} must equal the checksums in
  \code{SHA256SUMS} of this deposit (or, for a script adapted under point 4 to read the file, the checksum of the
  adapted script whose diff and check were posted; Appendix~\ref{app:p6}). The scripts compare these checksums only
  with the files on the computer that runs them, and \code{s110\_summary.py} takes the date of its 2027 February 28
  rule from that computer's clock. The public comparison is therefore the check that counts. A fallback primary needs
  the MAST search of Appendix~\ref{app:data}, and a variant declared not available needs its error output
  (Appendix~\ref{app:variants}). Appendix~\ref{app:public} lists every rule that only this public comparison checks.
\end{enumerate}
